% Prospective apparatus and protocol; scientific and publication acceptance remain open.
\section{Build a brain--body system we can test}

The proposed architecture becomes a first experiment only when its measurements
are specified. Begin with one hand and one movement that can be measured well.
Record the brain and muscle signals, capture the movement, and compare the
virtual hand's prediction from one declared input side with that independent
measurement. The error identifies what the next experiment must resolve.

The apparatus would use non-invasive EEG, surface EMG, a calibrated
motion-capture relay-composite, an open musculoskeletal-simulation physics-engine
with a lifelike view, and an open multimodal model trained on eligible neural
and muscular structure and activity data. The EEG/EMG observation pair and the
movement/physics pair would meet inside one conditional learner. Intent-related
brain signals and muscle activity would supply observations; measured movement
would check what the body actually did. Predictions from either end must answer
to observations outside the model. Form constrains what could happen; measured
function tests the prediction.

\subsection{One movement. Two directions.}

The forward test uses EEG and declared task context to predict supported muscle
activity and movement on withheld recordings. The reverse test uses muscle
activity and movement to predict recorded EEG features within that task. Compare
both predictions with independent measurements. Agreement between the models
cannot substitute for that test.

Movement measures what the body did. It narrows the explanations of brain
signals that remain plausible without uniquely identifying intent or subjective
experience. Retain stimulus, task, behaviour and separately reported context.
Show uncertainty instead of displaying a generated neural pattern as the
observed mind.

Surface EMG records muscle electrical activity. A claimed incoming sensory or
feedback pathway needs evidence of its own. The virtual body cannot supply that
missing observation. Keep recordings, predictions and simulation identifiable
throughout the test.

\subsection{A bridge must account for the spine}

The spinal cord helps shape the response. Muscle output depends on sensory
input, movement context, local circuitry and descending input; some spinal
pathways can retain learned changes. Reproducing a specified spinal function
would therefore require more than forwarding its signal: a brain--body bridge
must account for the processing it aims to reproduce.

A changed response can reflect immediate reflex modulation, lasting spinal
plasticity or peripheral adaptation. Those explanations need separate tests.
A change during the task does not establish retained memory, and a retained
change does not by itself locate that memory. Human reflex conditioning and
constrained mouse adaptation supply evidence for different claims.~\cite{thompson2009retained,thompson2013retained,lavaud2024}
The human studies do not establish brain-independent learning; the mouse
preparation does not demonstrate unrestricted walking or reaching. Our human
comparisons use abstracts and selected results, discussion or methods passages.
Relevant mouse-study passages were also examined in the author manuscript;
equivalence to the final published edition and the complete methods and
statistics have not been established here.

The periphery adds a different question: did the sensory signal itself change?
Passive proprioceptive training supplies a task-specific peripheral comparison;
its psychophysical and spindle-afferent experiments were separate.~\cite{ackerley2022retained}
Our passive-training comparison uses abstracts and selected results, discussion
or methods passages. None of these studies demonstrates universal skill
transfer or complex memory stored in a peripheral axon.

The next comparison must choose a function we can measure. Does the controller
respond appropriately to the supported context? Does an adaptation persist
through the declared test? Could peripheral or mechanical changes explain the
outcome? Keep recordings, estimates and simulated state in distinct visible
layers. Surface EMG cannot supply a missing sensory-nerve or spinal channel;
an explanatory cutaway cannot register recorded neural sites.

Externalization remains the research aim: can a specified spinal function be
reproduced outside its usual pathway? Replacing a computation, stimulating a
pathway, controlling a simulated limb and restoring biological function would
answer different questions. Each needs its own task- and site-specific evidence;
no programme spinal intervention or qualified spinal-writing interface is reported. Anatomical access does not establish
neural function, increased excitability does not establish useful plasticity,
and no fixed percentage of a complex skill can be allocated to brain, spine
or periphery from these proposals.

\subsection{Reading a signal and writing an output}

The first prototype would test the recording and prediction route with EEG,
EMG and independent movement measurements. To deliver feedback or act on a
body, it would also need the separate sensory or action route specified in
Section~\ref{b2-output-interfaces}. That method distinguishes indirect sensory
feedback from direct neural writing and virtual control from robotic or later
biological actuation, with the device, measurement, safety and clinical evidence
each would require.

Physical sensory writers, transducers and body actuators remain excluded from
this acquisition and virtual-work plan, unqualified and unpriced. Their absence
defines what this first prototype could test. The consolidated interface method
also separates those later physical procedures from the current observations: no stimulation, participant acquisition, biological actuation or qualified physical-interface training is reported.

\subsection{The first prototype has to earn the next one}

The lean tier tests whether the selected signals and calibrated movement
reference can be synchronized and used for a bounded task. The serious-prototype tier
improves acquisition, independent references and source-specific bidirectional
evaluation. Expand cohorts, tasks, instrumentation and compute in the research-lab
tier only when earlier evidence justifies the next step.

Equipment, calibration, accessories, storage, integration, independent references, safety and consent work, and contingency would all affect the cost. No complete, dated cost estimate has been established for these tiers. The lean, serious-prototype and research-lab scopes therefore specify different experiments, not demonstrated price points or a complete neural reconstruction. The unfulfilled currency, vendor and quotation requirements remain in the budget-preparation provenance record.

\subsection{Let the model earn the claim}

Which open recordings can support the next model comparison? Their rights,
version, coverage, clock, units and missing modalities determine the answer.
Training would proceed only under an independently admitted protocol, with test
participants, sessions and tasks kept outside model selection. Simple baselines
would give the added model a comparison to beat. Examining artifact leakage,
absent inputs, timing errors, physically invalid solutions and uncertainty would
keep that comparison answerable to its actual conditions. Earlier failures and
reserved recordings remain constraints on it. The new ambition must answer to
the evidence already kept.

Begin with constrained virtual hand or foot control. Robotics, remote operation,
unfamiliar effectors and extra limbs remain later research directions, each
requiring separate physical, latency and safety evidence. A successful movement
prediction supports a movement claim. Biological and perceptual claims require
the corresponding broader evidence.

Rights-cleared images and lifelike interactive models would make the apparatus and its predictions inspectable. The comparison should identify the observation, the prediction and the error, with the conditions that limit each claim. A convincing image cannot supply a missing measurement or demonstrate biological function. The next experiment must make those distinctions visible enough to test.